\documentclass{labreport}

% scripts/build_report.py が spec/REPORT_META.yaml から再生成する。
\reportcourse{科目名を指定}
\reporttitle{実験題目を指定}
\reportsubtitle{}
\reportfield{所属}{指定待ち}
\reportfield{学生番号}{指定待ち}
\reportfield{氏名}{指定待ち}
\reportfield{実験日}{指定待ち}
\reportfield{提出日}{指定待ち}



% 本文はユーザー原文から作成したファイルだけを読み込む。
% 未提出時は骨格PDFに（文章N）を表示する。
\IfFileExists{../build/report-mode.tex}{\input{../build/report-mode.tex}}{\def\reportskeleton{0}}
\newcommand{\slotplaceholder}[1]{%
  \par\medskip\noindent
  \fbox{\parbox{0.94\linewidth}{\centering\textbf{（#1）}}}%
  \par\medskip
}
\newcommand{\proseslot}[2]{%
  \ifnum\reportskeleton=1
    \slotplaceholder{#2}%
  \else
    \IfFileExists{../prose/typeset/#1.tex}{\input{../prose/typeset/#1.tex}}{\slotplaceholder{#2}}%
  \fi
}

\begin{document}
\makelabcover

% 表紙の次は、タイトル・氏名・概要・キーワードを再掲せず本文へ入る。
\section{目的}

TODO: 実習書の要求に合わせて目的を記す．章名と構成は実習書を優先する．


\section{実験方法}

\proseslot{P-002}{文章2}

\section{結果}

\proseslot{P-003}{文章3}

\section{考察}

\proseslot{P-004}{文章4}

\section{結論}

\proseslot{P-005}{文章5}

\begin{thebibliography}{9}
% 本文で実際に参照し、SOURCE_AUDIT.md で verified の文献だけを書く。
\end{thebibliography}



\end{document}
